\chapter{Conformite a l'Enonce}

\section{Table de conformite}
\begin{table}[H]
  \centering
  \caption{Conformite entre l'enonce et les realisations}
  \begin{tabular}{p{4.0cm}p{2.2cm}p{7.2cm}}
    \toprule
    Exigence de l'enonce & Etat & Preuve / element du projet \\
    \midrule
    Extraction de narratives depuis le corps Wikipedia & Fait & \texttt{scripts/extract\_timelines.py}, fichiers \texttt{*.timeline.json} \\
    Gestion du contexte long (chunking) & Fait & \texttt{scripts/chunk\_texts.py} \\
    Traitement multilingue FR/EN & Fait & pipeline retenu: \texttt{scripts/translate\_fr\_timelines\_to\_en.py} puis paires \texttt{en <-> fr\_en} \\
    Fine-tuning sentence-encoder (AutoTrain) & Fait & modeles \texttt{allMiniLM} et \texttt{MPNET} fine-tunes sur \texttt{autotrain\_pair\_score\_en\_fr\_en\_hard\_final} \\
    Interface Spaces avec score + retries & Fait & \texttt{app.py}, \texttt{scripts/generate\_analogues.py} \\
    Comparaison de 2 modeles de tailles differentes sur la precision des faits extraits & Partiel & Non finalise dans cette version; priorisation des autres blocs techniques avec l'encadrant \\
    Collecte via dataset HF \texttt{wikimedia/wikipedia} + API MediaWiki & Partiel & Collecte realisee via API MediaWiki; non-integration du dataset HF dans cette version \\
    Categorie films/livres/personnalites & Partiel & Categorie retenue: personnalites (acteurs/actrices) \\
    Dataset de paires narratives (livrable) & Fait & jeux \texttt{jsonl} de paires (train/validation/test) \\
    Branche fusion semantique + finalisation EN unique & Exploratoire & \texttt{scripts/merge\_multilingual\_timelines.py}, \texttt{scripts/finalize\_monolingual\_en\_timelines.py} \\
    \bottomrule
  \end{tabular}
\end{table}

\section{Positionnement de cette version}
Le projet a ete cadre pour livrer un pipeline stable de bout en bout sur une categorie
homogene (acteurs/actrices), avec une evaluation quantitative solide de la similarite
narrative et un demonstrateur fonctionnel.

La comparaison inter-LLM sur la precision des faits extraits n'a pas ete finalisee dans
cette iteration. Le choix de priorisation a ete de consolider l'extraction,
la traduction FR->EN pour les paires cross-lingues, le fine-tuning et l'integration applicative.

\section{Justification de la collecte des donnees}
La collecte a ete implementee via l'API MediaWiki, choisie pour la simplicite d'integration,
la stabilite operationnelle et le controle direct sur les pages/langues ciblees.
L'integration explicite du dataset \texttt{wikimedia/wikipedia} reste une extension
possible d'industrialisation.
